\section{Equivalent growth conditions for irreducible positive maps}


\begin{theorem}[Exponential positivity for irreducible positive maps]
    \label{thm:exp_pd_irred}
    \lean{exp_posDef_of_growth, exp_posDef_of_irreducible,
        exp_posDef_of_irreducible_cp}
    \leanok
    \uses{def:positive_map, def:irreducible_cp}
    Let $E$ be an irreducible positive map on $\MN{D}$,
    let $A \ge 0$ be nonzero, and let $t > 0$.
    Then
    \begin{align}
        \exp(tE)(A)
         & = \sum_{k=0}^{\infty}\frac{t^k}{k!}E^k(A)>0.
        \label{eq:qpf_exp_positive}
    \end{align}
    This is the forward implication in
    \cite[Theorem~6.2(3)]{Wolf2012Quantum}. Complete positivity is not used;
    the former completely positive declaration is retained as a direct
    specialization.
\end{theorem}

\begin{proof}
    \leanok
    \uses{thm:exp_truncation_pd_irred}
    By Theorem~\ref{thm:exp_truncation_pd_irred}, the first $D$ terms in
    (\ref{eq:qpf_exp_positive}) already form a positive definite matrix.
    Every remaining term is positive semidefinite, so the full exponential
    series is the sum of a positive definite matrix and a positive
    semidefinite tail.
\end{proof}

\begin{theorem}[Exponential characterization of irreducibility]
    \label{thm:wolf_6_2_item_3}
    \lean{irreducible_of_exp_posDef_forall,
        irreducible_iff_exp_posDef_forall_of_positive,
        irreducible_iff_exp_posDef_forall}
    \leanok
    \uses{def:positive_map, def:irreducible_cp}
    Let $E$ be a positive map on $\MN{D}$. Then $E$ is
    irreducible if and only if, for every $t>0$ and every nonzero
    $A\geq0$, one has $\exp(tE)(A)>0$.
    This is the equivalence in
    \cite[Theorem~6.2(3)]{Wolf2012Quantum};
    Theorem~\ref{thm:exp_pd_irred} is its forward implication.
\end{theorem}

\begin{proof}\leanok
    \uses{thm:exp_pd_irred}
    The forward implication is Theorem~\ref{thm:exp_pd_irred}. Conversely,
    let $P$ be an invariant orthogonal projection. Invariance gives
    $E(P\MN{D}P)\subseteq P\MN{D}P$, so induction yields
    $E^k(P)\in P\MN{D}P$ for every $k\geq0$. Therefore every term, and hence
    the convergent series
    \begin{align}
        \exp(tE)(P)
         & =\sum_{k=0}^{\infty}\frac{t^k}{k!}E^k(P),
        \notag
    \end{align}
    lies in the same corner. If $P\neq0$, then at $t=1$ the assumed positive
    definiteness is impossible unless $P=\Id$, because a positive definite
    matrix cannot be supported on a proper corner. Thus every invariant
    projection is $0$ or $\Id$, so $E$ is irreducible. The former completely
    positive declaration is retained as a direct specialization.
\end{proof}

\begin{theorem}[Irreducible positive maps]
    \label{thm:wolf_6_2}
    \lean{wolf_theorem_6_2_tfae,
        irreducible_of_orthogonal_trace_pos_forall}
    \leanok
    \uses{def:positive_map, def:irreducible_cp, thm:growth_pd_irred,
        thm:exp_pd_irred, thm:orthogonal_trace_cond}
    Let $E:\MN{D}\to\MN{D}$ be a positive linear map. The following
    conditions are equivalent:
    \begin{enumerate}
        \item If a Hermitian projection $P$ satisfies
              $E(P\MN{D}P)\subseteq P\MN{D}P$, then
              $P\in\{0,\Id\}$.
        \item For every nonzero $A\geq0$,
              $(\Id+E)^{D-1}(A)>0$.
        \item For every nonzero $A\geq0$ and every $t>0$,
              $\exp(tE)(A)>0$.
        \item For every orthogonal pair of nonzero matrices $A,B\geq0$,
              there is an integer $t\in\{1,\ldots,D-1\}$ such that
              $\tr(BE^t(A))>0$.
    \end{enumerate}
    This is \cite[Theorem~6.2]{Wolf2012Quantum}, with the source's positive-map
    hypothesis and without a Kraus or complete-positivity assumption.
\end{theorem}

\begin{proof}
    \leanok
    \uses{thm:growth_pd_irred, thm:exp_pd_irred,
        thm:orthogonal_trace_cond}
    The implication (1)$\Rightarrow$(2) is the support-projection and kernel
    descent argument of Theorem~\ref{thm:growth_pd_irred}.
    For (2)$\Rightarrow$(3), suppose that a vector lies in the kernel of the
    first $D$ terms of the exponential series. Positivity of the summands
    forces $E^k(A)$ to annihilate that vector for $0\leq k\leq D-1$.
    The binomial expansion of $(\Id+E)^{D-1}(A)$ then annihilates the same
    vector, contradicting (2). For (3)$\Rightarrow$(1), an invariant corner contains
    every $E^k(P)$ and hence $\exp(E)(P)$; positive definiteness then excludes
    a proper nonzero $P$. The implication (2)$\Rightarrow$(4) is
    Theorem~\ref{thm:orthogonal_trace_cond}. Finally, if (4) held while a
    proper nonzero projection $P$ defined an invariant corner, then every
    $E^t(P)$ would remain supported on $P$, and
    $\tr((\Id-P)E^t(P))=0$ for every $t$. This contradicts (4) with
    $A=P$ and $B=\Id-P$.
\end{proof}

\section{Ergodicity of irreducible channels}

\begin{theorem}[Unique density-matrix fixed point for an irreducible channel]
    \label{thm:channel_density_fp_irred}
    \lean{IsChannel.exists_unique_density_fixedPoint_of_irreducible}
    \leanok
    \uses{def:channel, def:irreducible_cp}
    Let $E$ be an irreducible quantum channel on $\MN{D}$ with $D > 0$.
    Then there exists a unique density matrix $\sigma$ such that
    $E(\sigma) = \sigma$.
    The fixed point $\sigma$ is positive definite.
    This is the fixed-point conclusion in the forward direction of
    \cite[Corollary~6.3]{Wolf2012Quantum}, specialized to completely
    positive trace-preserving maps.
\end{theorem}

\begin{proof}
    \leanok
    \uses{lem:cesaro_mean_density, lem:cesaro_subseq_fixed,
        thm:psd_eigenvector_pd_irred, thm:psd_eigenvector_unique_irred}
    Every Ces\`aro mean (\ref{eq:channel_cesaro_mean}) of a density matrix
    stays inside the compact convex set of density matrices, so a subsequence
    converges to a density matrix~$\sigma$.  The telescope identity
    (\ref{eq:channel_cesaro_telescope}) gives
    $E(\sigma)=\sigma$.  Theorem~\ref{thm:psd_eigenvector_pd_irred}
    makes $\sigma$ positive definite.  If $\tau$ is another density-matrix
    fixed point, Theorem~\ref{thm:psd_eigenvector_unique_irred} gives
    $\tau=c\sigma$.  Taking traces yields $1=c$, hence $\tau=\sigma$.
\end{proof}

\begin{theorem}[Ces\`aro convergence for irreducible channels]
    \label{thm:channel_cesaro_irred}
    \lean{IsChannel.cesaroMean_tendsto_of_irreducible}
    \leanok
    \uses{def:channel, def:irreducible_cp, def:cesaro_mean}
    Let $E$ be an irreducible quantum channel on $\MN{D}$ with $D > 0$,
    and let $\rho$ be any density matrix.
    Then the Ces\`aro means (\ref{eq:channel_cesaro_mean}) converge to the
    unique positive definite density-matrix fixed point of~$E$.
    This is the Ces\`aro-convergence conclusion in the forward direction of
    \cite[Corollary~6.3]{Wolf2012Quantum}, specialized to quantum channels.
\end{theorem}

\begin{proof}
    \leanok
    \uses{thm:channel_density_fp_irred}
    Every subsequential limit of the Ces\`aro means is again a density-matrix
    fixed point.
    By Theorem~\ref{thm:channel_density_fp_irred}, there is only one such
    fixed point, so every convergent subsequence has the same limit.
    Compactness then forces the whole sequence to converge to that limit.
\end{proof}

\section{Spectral characterization of irreducibility for positive maps}

\begin{definition}[Wolf spectral properties]
    \label{def:wolf_spectral_properties}
    \lean{HasWolfSpectralProperties}
    \leanok
    A linear map $T$ on $\MN{D}$ has the \emph{Wolf spectral properties} if
    there are a real number $r>0$ and matrices $X,Y>0$ such that
    \begin{align}
        \varrho(T)&=r, \notag\\
        T(X)&=rX, \notag\\
        T^*(Y)&=rY,
        \label{eq:qpf_wolf_spectral_eigenvectors}
    \end{align}
    and the ordinary complex eigenspace $\ker(T-r\Id)$ has dimension one.
    Here $T^*$ is the adjoint for the bilinear trace pairing.  The last clause
    is geometric nondegeneracy; it is not restricted to positive semidefinite
    eigenvectors and does not assert algebraic simplicity.
\end{definition}

\begin{theorem}[Irreducibility gives the Wolf spectral properties]
    \label{thm:wolf_spectral_properties_irred_positive}
    \lean{hasWolfSpectralProperties_of_irreducible_positive}
    \leanok
    \uses{def:wolf_spectral_properties, def:positive_map,
        def:irreducible_cp}
    Let $D>0$, and let $T:\MN{D}\to\MN{D}$ be a nonzero irreducible
    positive map.  Then $T$ has the Wolf spectral properties of
    Definition~\ref{def:wolf_spectral_properties}.
\end{theorem}

\begin{proof}
    \leanok
    \uses{thm:wolf_irreducible_positive_spectral_radius,
        thm:trace_adjoint_irreducible_positive}
    The positive-map form of Wolf Theorem~6.3 gives a positive-definite right
    eigenvector at the spectral radius and proves that its ordinary eigenspace
    is one-dimensional.  Applying the trace-adjoint consequence of the same
    theorem gives a positive-definite left eigenvector at the same eigenvalue.
\end{proof}

\begin{theorem}[Wolf spectral properties imply irreducibility]
    \label{thm:irred_from_wolf_spectral_properties}
    \lean{isIrreducibleMap_of_hasWolfSpectralProperties}
    \leanok
    \uses{def:wolf_spectral_properties, def:positive_map,
        def:irreducible_cp}
    Let $D>0$, and let $T:\MN{D}\to\MN{D}$ be positive.  If $T$ has the
    Wolf spectral properties, then $T$ is irreducible.
\end{theorem}

\begin{proof}
    \leanok
    \uses{lem:similarity_irred, lem:cesaro_subseq_fixed}
    Let $r,X,Y$ be the data in
    (\ref{eq:qpf_wolf_spectral_eigenvectors}), put $S=\sqrt{Y}$, and define
    exactly as in Wolf Equation~(6.34)
    \begin{align}
        T'(Z)
          &=r^{-1}S\,T(S^{-1}ZS^{-1})S.
        \label{eq:qpf_wolf_spectral_tp_gauge}
    \end{align}
    The map $T'$ is positive.  Trace duality and $T^*(Y)=rY$ give
    \begin{align}
        \tr(T'(Z))
          &=r^{-1}\tr(YT(S^{-1}ZS^{-1}))
           =\tr(Z),
        \notag
    \end{align}
    so $T'$ is trace preserving, and $SXS>0$ is fixed by $T'$.

    If $T'$ preserved a nonzero proper corner, positive trace-preserving
    Ces\`aro averaging inside that corner would give a stationary density
    matrix there.  Conjugating it by $S^{-1}$ produces an ordinary
    $r$-eigenvector of $T$.  Eigenspace dimension one makes the corner state a
    scalar multiple of $SXS$, which is impossible because $SXS$ has full
    support.  Hence $T'$ is irreducible, and similarity invariance gives
    irreducibility of $T$.
\end{proof}

\begin{theorem}[Wolf Theorem 6.4]
    \label{thm:wolf_theorem_6_4}
    \lean{wolf_theorem_6_4}
    \leanok
    \uses{def:wolf_spectral_properties, def:positive_map,
        def:irreducible_cp}
    Let $D>0$, and let $T:\MN{D}\to\MN{D}$ be a nonzero positive map.
    Then $T$ is irreducible if and only if it has the Wolf spectral properties.
    This is \cite[Theorem~6.4]{Wolf2012Quantum}, local source lines~698--721,
    with one necessary boundary correction: on $\MN{1}$ the zero map is
    irreducible but has spectral radius zero, so it cannot satisfy the printed
    strict-positive condition $T(X)=rX>0$.
\end{theorem}

\begin{proof}
    \leanok
    \uses{thm:wolf_spectral_properties_irred_positive,
        thm:irred_from_wolf_spectral_properties}
    Combine the preceding two theorems.
\end{proof}

\subsection{The earlier completely positive specialization}

\begin{definition}[Restricted CP spectral properties]
    \label{def:has_spectral_properties}
    \lean{HasSpectralProperties}
    \leanok
    \uses{def:cp_map}
    A completely positive map $E$ on $\MN{D}$ has the \emph{restricted CP
        spectral properties} used here if there exist a positive real number $r$, a positive
    definite right eigenvector $\rho$, and a positive definite left eigenvector
    $\sigma$ for the adjoint map such that
    \begin{align}
        E(\rho)
         & =r\rho, \notag \\
        E^\dagger(\sigma)
         & =r\sigma,
        \label{eq:qpf_spectral_eigenvectors}
    \end{align}
    every positive semidefinite right eigenvector for eigenvalue $r$ is a
    scalar multiple of $\rho$, and the spectral radius of $E$ is equal to~$r$.
    The uniqueness clause concerns only positive semidefinite Perron
    eigenvectors. Unlike the full nondegenerate-eigenspace statement in
    \cite[Theorem~6.4]{Wolf2012Quantum}, it does not assert that every complex
    eigenvector at $r$ is proportional to $\rho$.
\end{definition}

\begin{theorem}[Irreducibility gives the restricted CP spectral properties]
    \label{thm:has_spectral_properties_irred}
    \lean{hasSpectralProperties_of_irreducible_cp}
    \leanok
    \uses{def:has_spectral_properties, def:irreducible_cp}
    Every nonzero irreducible completely positive map on $\MN{D}$ has the
    restricted CP spectral properties of Definition~\ref{def:has_spectral_properties}.
\end{theorem}

\begin{proof}
    \leanok
    \uses{thm:pf_eigenvector_existence, thm:psd_eigenvector_unique_irred,
        thm:spectral_radius_pf_irred}
    Combine Perron--Frobenius existence for the map and its adjoint with the
    uniqueness theorem for positive semidefinite Perron eigenvectors and the
    spectral-radius identity of
    Theorem~\ref{thm:spectral_radius_pf_irred}.
\end{proof}

\begin{theorem}[Restricted CP spectral properties imply irreducibility]
    \label{thm:irred_from_spectral_properties}
    \lean{isIrreducibleMap_of_hasSpectralProperties}
    \leanok
    \uses{def:has_spectral_properties}
    If a completely positive map on $\MN{D}$ has the restricted CP spectral properties of
    Definition~\ref{def:has_spectral_properties}, then it is irreducible.
\end{theorem}

\begin{proof}
    \leanok
    \uses{thm:kraus_tp_gauge_from_adjoint_fixed, lem:cesaro_subseq_fixed}
    Gauge by the positive definite left eigenvector in
    (\ref{eq:qpf_spectral_eigenvectors}) and rescale by the Perron eigenvalue
    to obtain a trace-preserving map with a positive definite fixed point.
    A nontrivial invariant projection would then produce, by Ces\`aro averaging
    in its corner algebra, a second positive semidefinite fixed point not
    proportional to the Perron vector, contradicting the uniqueness clause in
    Definition~\ref{def:has_spectral_properties}.
\end{proof}

\begin{theorem}[Irreducibility iff the restricted CP spectral properties hold]
    \label{thm:irred_iff_spectral_properties}
    \lean{isIrreducibleMap_iff_spectral_properties}
    \leanok
    \uses{def:has_spectral_properties, def:irreducible_cp}
    Let $E$ be a nonzero completely positive map on $\MN{D}$.
    Then $E$ is irreducible if and only if it has the restricted CP spectral properties of
    Definition~\ref{def:has_spectral_properties}. This is a CP-map variant of
    \cite[Theorem~6.4]{Wolf2012Quantum}: uniqueness is required only among
    positive semidefinite Perron eigenvectors, not on the full eigenspace.
    Those restricted properties are nevertheless sufficient for the
    irreducibility equivalence.
\end{theorem}

\begin{proof}
    \leanok
    \uses{thm:has_spectral_properties_irred, thm:irred_from_spectral_properties}
    Combine Theorems~\ref{thm:has_spectral_properties_irred}
    and~\ref{thm:irred_from_spectral_properties}.
\end{proof}
